\section{Introduction}

\myFloat{primitives}{t}{0.9}

The capability of Large Language Models (LLMs) to process extensive contexts has revolutionized applications ranging from long-document summarization \cite{longbenchv2}, repository-level code analysis \cite{coderepoqa}, to long-term agentic planning \cite{planandact}. However, this capability comes at a steep serving cost. Long-context LLM inference is bottlenecked by both the quadratic attention complexity and the linear growth of the Key-Value (KV) cache. In long-context scenarios, the sheer volume of tokens not only incurs massive computational overhead but also exhausts limited GPU memory. Consequently, the latency and resource consumption increase significantly with the context length, rendering long-context LLM inference prohibitively expensive.

To mitigate this bottleneck, prior research has focused on two management strategies: \textbf{KV Selection} and \textbf{KV Eviction}. Selection-based methods, such as Quest \cite{quest} and NSA \cite{nsa}, selectively attend to relevant KVs at runtime yet retain the complete state in memory. Eviction-based methods, such as SnapKV \cite{snapkv} and R-KV \cite{rkv}, retrospectively prune the cache by evicting unneeded tokens after they have accumulated in memory. While effective, these paradigms share a common, fundamental inefficiency: \textit{every generated token is indiscriminately admitted to the growing cache state}, where the system blindly consumes memory resources to accumulate redundant KV states, only to later skip or evict them.

In this paper, we argue that an efficient long-context system requires a third, missing primitive: \textbf{KV Admission}. Instead of retrieving information at \textit{read-time} (Selection) or managing memory \textit{post-write} (Eviction), Admission operates \textit{pre-write}, acting as a gatekeeper that predicts the future utility of a token before admitting it to the KV cache.

To realize this, we introduce \textbf{Write-Gated KV (WG-KV)}, a novel, learnable KV Admission mechanism. Unlike heuristic approaches, WG-KV employs a lightweight, end-to-end differentiable predictor that evaluates token utility \textit{before} cache entry. During prefilling, this predictor constructs a vertical-slash attention mask, ensuring high-utility tokens are visible to the entire sequence while others are accessible only locally. Simultaneously, it routes generated KV pairs to distinct memory regions: tokens predicted as high-utility are admitted to a persistent global cache, while others are confined to a transient local cache. This selective strategy seamlessly extends to the decoding phase, where we maintain this dual-cache structure to filter redundant states upon KV generation.

Our approach yields substantial efficiency gains through a \textbf{hardware-aware implementation} that effectively translates theoretical sparsity into practical wall-clock speedups. We evaluate WG-KV on Llama 3 \cite{llama3} and Qwen3 \cite{qwen3} models, demonstrating that WG-KV reduces memory usage by 36--69\% and delivers 2.56--4.17x prefill and 1.59--2.63x decode speedups, while maintaining near-lossless task accuracy across diverse long-context benchmarks. By \textit{learning what to write}, we enable LLMs to sustain longer contexts with significantly fewer resources.